\chapter{Ionosphere, Plasmasphere, and Ray Tracing}
\label{chap:ionosphere}

This chapter specifies the ionospheric and plasmaspheric group-delay model that
\lupnt{} applies to terrestrial GNSS signals received by cislunar and lunar
users.  A lunar receiver observes the terrestrial GNSS constellation through the
side lobes of the transmit antennas, so the useful links graze the Earth: the
line of sight passes almost tangentially through the topside ionosphere and then
runs for tens of Earth radii through the plasmasphere.  Neither the thin-shell
mapping of the classical GNSS literature nor the Klobuchar broadcast model is applicable
to such geometry, because both assume a single, steeply inclined crossing of a
compressed ionospheric shell \citep{misra2011gnss}.  \lupnt{} instead integrates
the electron density along the numerically ray-traced, bent signal path.

The model is specified in \citep{iiyama2026gcpm}, with the earlier conference
treatment in \citep{iiyama2025iono} and the full characterization in
\citep[Ch.~5]{iiyama2026thesis}; equation numbers of the form~(5.x) quoted below
refer to that chapter.  The implementation lives in the \code{pecsim} namespace
under \file{cpp/lupnt/environment/plasma/}:

\begin{itemize}[nosep]
  \item \file{tec/raytrace.h}, \file{tec/raytrace.cc} --- the tracer proper
        (\cpp{pecsim::trace\_ray}, \cpp{RayTraceConfig}, \cpp{PathProfile});
  \item \file{gcpm/gcpm_interface.h}, \file{gcpm/gcpm_interface.cc} --- the
        Global Core Plasma Model (GCPM) v2.4 electron density
        \citep{gallagher2000gcpm};
  \item \file{gcpm/iri_interface.h}, \file{gcpm/iri_interface.cc} --- the
        International Reference Ionosphere (IRI) submodel
        \citep{bilitza2022iri} and its solar-index options;
  \item \file{igrf/igrf_interface.h}, \file{igrf/igrf_interface.cc} --- the
        IGRF-14 geomagnetic field \citep{alken2021igrf};
  \item \file{tec/iono_model.h} --- the \cpp{IonoModel} backend selector;
  \item \file{tec/raytrace_wrapper.h}, \file{tec/raytrace_wrapper.cc} --- the
        batch driver used by the precompute path;
  \item \file{tec/neldermead.h}, \file{tec/neldermead.cc} --- the simplex solver
        used by the receiver-side convergence iteration.
\end{itemize}

The delay is consumed by the GNSS measurement model through
\cpp{GNSSMeasurements::\allowbreak ComputeIonosphere\allowbreak PlasmaRayTraceDelay}
(\file{cpp/lupnt/measurements/gnss_measurement.cc}), where it appears as
$\Delta\rho_\mathrm{plasma}$ in \cref{chap:gnss-measurements}.

\section{Scope, Units, and the Frame Contract}
\label{sec:iono-contract}

The tracer is a self-contained module with its own unit system, inherited from
the Fortran heritage of GCPM and IRI.  Every caller is responsible for
converting at the boundary.

\begin{lupntnote}
The entire plasma module is \textbf{km-based}.  The positions \code{x1},
\code{x2} passed to \cpp{trace\_ray} are in kilometres, the step size, gradient
step and cutoff radius are in kilometres, the Earth radius is
$R_E = \SI{6371.0}{\kilo\meter}$ and the speed of light is
$c = \SI{299792.458}{\kilo\meter\per\second}$
(\file{cpp/lupnt/environment/plasma/core/constants.h}).  Electron density is
evaluated internally in \si{\per\cubic\centi\meter} by GCPM and converted to
\si{\per\cubic\meter} before entering the delay integrals; one TEC unit is
$\SI{1}{\TECU} = \num{1e16}$ electrons per \si{\square\meter}.  Delays are
returned in \textbf{metres}.  SI callers such as \cpp{GNSSMeasurements} divide
their metre positions by $1000$ before calling \cpp{trace\_ray}, and use the
returned metres directly.
\end{lupntnote}

The tracer works in a geocentric, Earth-fixed Cartesian frame.  The GNSS
coupling converts the receiver and transmitter states from \code{options.frame}
into \code{raytrace\_frame} (default \code{Frame::ECEF}) and supplies an epoch in
\code{raytrace\_epoch\_scale} (default \code{Time::UTC}) expressed as seconds
from J2000:
%
\begin{equation}
  \label{eq:iono-inputs}
  \vec{x}_1 = \vec{r}_\mathrm{tx}^{\,\mathrm{ecef}}\,[\si{\kilo\meter}],
  \qquad
  \vec{x}_2 = \vec{r}_\mathrm{rx}^{\,\mathrm{ecef}}\,[\si{\kilo\meter}],
  \qquad
  t = t_\mathrm{rx}^{\UTC}\ \text{(s since J2000)} .
\end{equation}
%
The unit conversion applied to every density sample is
%
\begin{equation}
  \label{eq:iono-ne-units}
  n_e\,[\si{\per\cubic\meter}] = 10^{6}\, n_e\,[\si{\per\cubic\centi\meter}] ,
  \qquad
  \mathrm{TEC}\,[\si{\TECU}] = 10^{-16} \int_\Gamma n_e\,\dd s\,[\si{\per\square\meter}] .
\end{equation}
%
Because the arc length $s$ is carried in kilometres while $n_e$ is in
\si{\per\cubic\meter}, every path integrand in \cref{sec:iono-refractive}
carries an explicit factor $1000$ in the code.

\Cref{tab:iono-symbols} fixes the symbols used throughout the chapter.

\begin{table}[htbp]
  \centering
  \caption{Symbols used in the ionosphere/plasmasphere model.}
  \label{tab:iono-symbols}
  \small
  \begin{tabularx}{\textwidth}{llL}
    \toprule
    Symbol & Unit & Meaning \\
    \midrule
    $n_e$              & \si{\per\cubic\meter}   & Electron number density (GCPM output, converted from \si{\per\cubic\centi\meter}) \\
    $f$                & \si{\hertz}             & Carrier frequency of the link (\code{config.freq\_Hz}) \\
    $f_p$              & \si{\hertz}             & Electron plasma frequency, \cref{eq:iono-fp} \\
    $f_g$              & \si{\hertz}             & Electron gyro (cyclotron) frequency, \cref{eq:iono-fp} \\
    $\vec{B}$, $B$     & \si{\nano\tesla}        & Geomagnetic field vector and its magnitude (IGRF-14) \\
    $\theta$           & rad                     & Angle between $\vec{B}$ and the local ray direction $\uvec{s}$ \\
    $n$, $n_\mathrm{gr}$ & ---                   & Phase and group refractive indices \\
    $\Gamma$           & ---                     & The (bent) ray path from transmitter to receiver \\
    $s$                & \si{\kilo\meter}        & Arc length along $\Gamma$ \\
    $\uvec{s}$         & ---                     & Unit tangent to $\Gamma$ \\
    $\rho$             & \si{\kilo\meter}        & Straight-line (chord) distance $\lVert\vec{x}_2-\vec{x}_1\rVert$ \\
    $s_f$              & \si{\kilo\meter}        & Total traced arc length (\code{PathProfile::sf}) \\
    $r$                & \si{\kilo\meter} or $R_E$ & Geocentric radius (in $R_E$ when passed to GCPM) \\
    $K_p$              & ---                     & Planetary geomagnetic activity index (\code{config.kp}) \\
    $R_{12}$           & ---                     & 12-month running mean sunspot number (\code{config.rz12}) \\
    $F_{10.7}$         & sfu                     & Solar radio flux at \SI{10.7}{\centi\meter} \\
    $IG_{12}$          & ---                     & Ionospheric (effective) global index used by IRI for $f_oF_2$ \\
    $\mathrm{MLT}$     & h                       & Solar-magnetic local time \\
    $\phi_\mathrm{SM}$ & rad                     & Solar-magnetic latitude \\
    \bottomrule
  \end{tabularx}
\end{table}

\section{Solar-Magnetic Coordinates}
\label{sec:iono-sm}

GCPM is parameterized in solar-magnetic (SM) coordinates: $\uvec{z}_\mathrm{SM}$
is aligned with the geomagnetic dipole axis, $\uvec{x}_\mathrm{SM}$ lies in the
plane spanned by the dipole axis and the Earth--Sun line (tilted towards the
Sun), and $\uvec{y}_\mathrm{SM}$ completes the right-handed triad.  SM is the
natural frame for the plasmasphere because the cold-plasma distribution is
organized by dipole $L$-shells and by local time relative to the Sun.

The conversion is a chain of four elementary rotations
(\file{cpp/lupnt/environment/plasma/gcpm/conversions.cc}), applied to the
Earth-fixed (geographic, \code{GEO} $\equiv$ ITRF/ECEF) input position:
%
\begin{equation}
  \label{eq:iono-sm-chain}
  \vec{x}_\mathrm{SM}
  = \mat{T}_4\,\mat{T}_3\,\mat{T}_2\,\mat{T}_1^{-1}\,\vec{x}_\mathrm{GEO},
  \qquad
  \underbrace{\mathrm{GEO}}_{\text{ITRF}}
  \xrightarrow{\ \mat{T}_1^{-1}\ } \underbrace{\mathrm{GEI}}_{\text{GCRF}}
  \xrightarrow{\ \mat{T}_2\ } \mathrm{GSE}
  \xrightarrow{\ \mat{T}_3\ } \mathrm{GSM}
  \xrightarrow{\ \mat{T}_4\ } \mathrm{SM} .
\end{equation}
%
With the passive elementary rotations
%
\begin{equation}
  \mat{R}_x(\alpha) =
  \begin{bmatrix} 1 & 0 & 0 \\ 0 & \cos\alpha & \sin\alpha \\ 0 & -\sin\alpha & \cos\alpha \end{bmatrix},
  \quad
  \mat{R}_y(\alpha) =
  \begin{bmatrix} \cos\alpha & 0 & \sin\alpha \\ 0 & 1 & 0 \\ -\sin\alpha & 0 & \cos\alpha \end{bmatrix},
  \quad
  \mat{R}_z(\alpha) =
  \begin{bmatrix} \cos\alpha & \sin\alpha & 0 \\ -\sin\alpha & \cos\alpha & 0 \\ 0 & 0 & 1 \end{bmatrix},
\end{equation}
%
the four stages are as follows.  Let $T_0 = (\mathrm{MJD} - 51544.5)/36525$ be
the time in Julian centuries from J2000 and let $\mathrm{ut}$ be the UT hour of
day.

\paragraph{Stage 1 (GEI $\leftrightarrow$ GEO).}  A rotation about the polar
axis by the Greenwich mean sidereal angle,
%
\begin{equation}
  \label{eq:iono-t1}
  \mat{T}_1 = \mat{R}_z(\vartheta),
  \qquad
  \vartheta = \left(100.461 + 36000.770\,T_0 + 15.04107\,\mathrm{ut}\right)^{\circ} .
\end{equation}
%
The tracer applies $\mat{T}_1^{-1} = \mat{R}_z(-\vartheta)$ to lift the ECEF
input into the quasi-inertial GEI frame.

\paragraph{Stage 2 (GEI $\to$ GSE).}  A rotation into the ecliptic and then to
the Sun direction,
%
\begin{equation}
  \label{eq:iono-t2}
  \mat{T}_2 = \mat{R}_z(\lambda_\odot)\,\mat{R}_x(\varepsilon),
\end{equation}
%
with obliquity, mean anomaly, mean longitude and apparent solar longitude
%
\begin{align}
  \varepsilon    &= \left(23.439 - 0.013\,T_0\right)^{\circ}, \\
  M              &= \left(357.528 + 35999.05\,T_0 + 0.04107\,\mathrm{ut}\right)^{\circ}, \\
  \gamma         &= \left(280.46 + 36000.772\,T_0 + 0.04107\,\mathrm{ut}\right)^{\circ}, \\
  \lambda_\odot  &= \gamma + \left(1.915 - 0.0048\,T_0\right)^{\circ}\sin M
                    + 0.020^{\circ}\sin 2M .
\end{align}

\paragraph{Stages 3 and 4 (GSE $\to$ GSM $\to$ SM).}  Let $\vec{q}$ be the
geomagnetic dipole unit vector expressed in GSE.  It is built from the IGRF
dipole pole in geographic coordinates and rotated forward through stages~1
and~2:
%
\begin{equation}
  \label{eq:iono-dipole}
  \phi_\mathrm{dip} = \left(78.8 + \num{4.283e-2}\,F\right)^{\circ},
  \qquad
  \lambda_\mathrm{dip} = \left(289.1 - \num{1.413e-2}\,F\right)^{\circ},
  \qquad
  F = \frac{\mathrm{MJD} - 46066}{365.25} ,
\end{equation}
%
so that $\vec{q}_\mathrm{GEO}$ is the unit vector at latitude
$\phi_\mathrm{dip}$ and longitude $\lambda_\mathrm{dip}$, and
$\vec{q} = \mat{T}_2\mat{T}_1^{-1}\vec{q}_\mathrm{GEO}$.  Then
%
\begin{equation}
  \label{eq:iono-t3t4}
  \mat{T}_3 = \mat{R}_x(\psi), \quad \psi = -\arctan\!\left(\frac{q_y}{q_z}\right),
  \qquad
  \mat{T}_4 = \mat{R}_y(\mu), \quad
  \mu = -\arctan\!\left(\frac{q_x}{\sqrt{q_y^2+q_z^2}}\right) .
\end{equation}
%
$\mat{T}_3$ rolls the frame about the Earth--Sun line until the dipole axis lies
in the $x$--$z$ plane (giving GSM); $\mat{T}_4$ pitches by the dipole tilt angle
$\mu$ until $\uvec{z}$ coincides with the dipole axis (giving SM).  When $q_z=0$
the code substitutes $\psi = -\tfrac{\pi}{2}\operatorname{sign}(q_y)$.

\paragraph{Spherical parameters passed to GCPM.}  After the rotation the
Cartesian SM position is converted to spherical coordinates and to magnetic
local time:
%
\begin{equation}
  \label{eq:iono-mlt}
  r = \frac{\lVert\vec{x}_\mathrm{SM}\rVert}{R_E},
  \qquad
  \phi_\mathrm{SM} = \arctan\!\left(\frac{z_\mathrm{SM}}{\sqrt{x_\mathrm{SM}^2 + y_\mathrm{SM}^2}}\right),
  \qquad
  \mathrm{MLT} = 12 + \frac{12}{\pi}\,\lambda_\mathrm{SM}\ \ [\mathrm{h}],
\end{equation}
%
where $\lambda_\mathrm{SM} = \atantwo(y_\mathrm{SM}, x_\mathrm{SM})$ is wrapped
into $[0,2\pi)$ and $\mathrm{MLT}$ is reduced modulo \SI{24}{\hour}.  Noon
($\mathrm{MLT}=12$) therefore corresponds to $\lambda_\mathrm{SM}=0$, i.e.\ the
sunward side of the SM frame.

\begin{lstlisting}[language=cpp, caption={Solar-magnetic transformation and GCPM
call.  Source: \texttt{cpp/lupnt/environment/plasma/tec/raytrace.cc} ::
\code{compute\_ne}.}, label={lst:iono-computene}]
// cpp/lupnt/environment/plasma/tec/raytrace.cc :: compute_ne
double compute_ne(double t_j2000, const Vec3d& pos_geo, RayTraceConfig config, bool debug) {
  double kp = config.kp;                       // Kp index for the ionosphere model
  if (kp < 0) throw std::runtime_error("Kp index is not set. Please set the Kp index.");
  if (get_iri_model() == IRIModel::NONE)
    throw std::runtime_error("IRI model is not set. Please set the IRI model.");

  double alatr = 0.0, along = 0.0, r = 0.0;
  std::array<double, 3> pos_gei_arr = {pos_geo[0] / RE, pos_geo[1] / RE, pos_geo[2] / RE};
  std::array<double, 3> pos_sm;

  double mjd = tj2000_to_mjd(t_j2000);
  DateTime datetime = mjd_to_datetime(mjd);
  std::array<int, 2> itime = datetime_to_itime(datetime);

  geo_to_sm(itime, pos_gei_arr, pos_sm);       // GEO -> GEI -> GSE -> GSM -> SM
  cart_to_pol(pos_sm, alatr, along, r);        // r in RE, alatr in rad, along in [0, 2pi)
  double amlt = 12.0 + along / AMLTRAD;        // magnetic local time [hours]
  if (amlt > 24.0) amlt -= 24.0;

  std::vector<double> out;
  if (config.use_fortran_gcpm) out = gcpm_v24_fortran(datetime, r, amlt, alatr, kp);
  else                         out = gcpm_v24(datetime, r, amlt, alatr, kp);
  return out[0];                               // total electron density [cm^-3]
}
\end{lstlisting}

\begin{lupntnote}
The local variable is named \code{pos\_gei\_arr}, but the routine invoked is
\cpp{geo\_to\_sm}: the argument is interpreted as an Earth-fixed (GEO/ITRF)
position, consistent with \code{raytrace\_frame} defaulting to
\code{Frame::ECEF}.  The name is a historical artefact and does not indicate a
frame change.  Note also that \cref{eq:iono-t1,eq:iono-t2} are the low-precision
sidereal and solar series used by the GCPM support code; they carry
milliarcsecond-class errors that are irrelevant at the resolution of an
empirical plasma model, but they are \emph{not} the IAU~2006/2000A chain used
elsewhere in \lupnt{} (\cref{chap:frame-conversions}).
\end{lupntnote}

\section{Electron-Density Environment: GCPM v2.4 over IRI}
\label{sec:iono-density}

The electron density is supplied by the Global Core Plasma Model (GCPM) v2.4
\citep{gallagher2000gcpm}, an empirical model of the cold plasma of the inner
magnetosphere.  GCPM stitches together region-specific submodels --- the
ionosphere (delegated to IRI \citep{bilitza2022iri}), the plasmasphere, the
plasmapause, the mid-latitude trough, and the polar cap --- using
hyperbolic-tangent blending functions so that the composite density field is
continuous and differentiable, which is a prerequisite for the gradient-driven
ray tracing of \cref{sec:iono-eikonal}.

As a mathematical contract the density along the path is
%
\begin{equation}
  \label{eq:iono-ne-model}
  n_e = n_e\big(r,\ \phi_\mathrm{SM},\ \mathrm{MLT};\ K_p,\ R_{12},\ t\big) ,
\end{equation}
%
with the spatial arguments defined by \cref{eq:iono-mlt} and the remaining
arguments defined below.

\subsection{Region blending}
\label{sec:iono-blending}

Let $L = r/\cos^2\phi_\mathrm{SM}$ be the dipole $L$-shell of the sample point
(clamped by $\cos^2\phi_\mathrm{SM}\ge 10^{-5}$), and let $L_\mathrm{crit}$ be the
$L$-shell of the poleward edge of the auroral oval, obtained by bilinear
interpolation of the empirical table $P_N$ / $P_S$ in $(\mathrm{MLT},K_p)$.  The
transition half-width is $\Delta L = 2.0$.  GCPM uses the switching function
%
\begin{equation}
  \label{eq:iono-switchon}
  \sigma(L; L_\mathrm{crit}, \Delta L)
  = \frac{1}{2}\left[1 + \tanh\!\left(\frac{3.4534\,(L - L_\mathrm{crit})}{\Delta L}\right)\right] ,
\end{equation}
%
which rises from $0$ to $1$ to within \SI{0.1}{\percent} as $L$ crosses
$L_\mathrm{crit} \pm \Delta L$, and blends the plasmasphere/trough branch with
the polar-cap branch:
%
\begin{equation}
  \label{eq:iono-blend}
  n_e =
  \begin{cases}
    n_e^\mathrm{ps/trough}, & L < L_\mathrm{crit} - \Delta L, \\[4pt]
    (1-\sigma)\,n_e^\mathrm{ps/trough} + \sigma\,n_e^\mathrm{cap},
      & L_\mathrm{crit} - \Delta L \le L \le L_\mathrm{crit} + \Delta L, \\[4pt]
    n_e^\mathrm{cap}, & L > L_\mathrm{crit} + \Delta L .
  \end{cases}
\end{equation}
%
Both branches call IRI below the F2 peak and bridge smoothly to the
plasmaspheric/polar profiles above it.  The same $\sigma$ controls the ion
composition: with $h = (r-1)R_E$ the altitude in \si{\kilo\meter},
%
\begin{align}
  \label{eq:iono-composition}
  \frac{n_{\mathrm{He}^+}}{n_{\mathrm{H}^+}}
    &= 10^{\,-1.541 - 0.176\,r + \num{8.557e-3}F_{10.7} - \num{1.458e-5}F_{10.7}^2}
       \left[1 - \sigma(L; L_\mathrm{crit}, \Delta L)\right], \\
  \alpha_{\mathrm{O}^+}
    &= \frac{0.995}{\left[1 + \left(\dfrac{h - 350}{531.25}\right)^{2}\right]^{3}} + 0.005 ,
\end{align}
%
from which the $\mathrm{H}^+$, $\mathrm{He}^+$ and $\mathrm{O}^+$ partial
densities are recovered.
Only the \emph{total} electron density (output slot~0) enters the delay
integrals; the ion partition is carried for diagnostics.

\begin{lstlisting}[language=cpp, caption={GCPM v2.4 region blending.  Source:
\texttt{cpp/lupnt/environment/plasma/gcpm/gcpm\_interface.cc} ::
\code{gcpm\_v24}.}, label={lst:iono-gcpm}]
// cpp/lupnt/environment/plasma/gcpm/gcpm_interface.cc :: gcpm_v24
IRIParams params = iri_sm(0.0, 0.0, 1.01, itime);   // f107 / rz12 / hmF2 from IRI
double f107 = params.f107, rz12 = params.rz12, hmf2_km = params.hmf2_km;

alcrit   = 1.0 / pow(cos(alatcritn * DEG2RAD), 2);  // L-shell of the auroral edge
tranlow  = alcrit - altrans;                        // altrans = 2.0
tranhigh = alcrit + altrans;

clat = pow(cos(alatr), 2);
if (clat < 1.0e-5) clat = 1.0e-5;
al      = r / clat;                                 // dipole L-shell
aheight = (r - 1.0) * re;                           // altitude [km]

double edensity = 0.0;
if (al < tranlow) {
  edensity = ne_iri_ps_trough(r, al, alatr, amlt, akp, itime);
} else if (al <= tranhigh) {
  double ps_edensity  = ne_iri_ps_trough(r, al, alatr, amlt, akp, itime);
  double cap_edensity = ne_iri_cap(r, alatr, amlt, itime);
  double switch_val   = switchon(al, alcrit, altrans);
  edensity = ps_edensity * (1.0 - switch_val) + cap_edensity * switch_val;
} else {
  edensity = ne_iri_cap(r, alatr, amlt, itime);
}
double den = edensity / 1.0e6;                      // -> particles per cm^3
\end{lstlisting}

\subsection{Model inputs}
\label{sec:iono-inputs}

\begin{description}[style=nextline, leftmargin=1.2em]
  \item[$K_p$ --- \code{config.kp}]
    The planetary geomagnetic activity index.  It sets the plasmapause radius:
    higher $K_p$ erodes the plasmasphere, contracting the plasmapause from
    $\sim$4--6\,$R_E$ under quiet conditions to $\sim$2--3\,$R_E$ under active
    conditions, and it also drives the auroral-edge table that fixes
    $L_\mathrm{crit}$ in \cref{eq:iono-switchon}.  \cpp{compute\_ne} throws when
    $K_p < 0$.  Passing $K_p<0$ directly to \cpp{gcpm\_v24} instead triggers a
    lookup of the historical index for the requested epoch.
  \item[$R_{12}$ --- \code{config.rz12}]
    The 12-month running mean sunspot number, which parameterizes the IRI
    submodel.  See \cref{sec:iono-indices}.
  \item[$t$ --- the epoch]
    Sets both the IRI solar-activity state and the SM frame orientation through
    \cref{eq:iono-t1,eq:iono-t2,eq:iono-t3t4}.
\end{description}

\begin{lupntnote}
\code{config.kp} must be non-negative and an IRI model must have been selected
with \cpp{set\_iri\_model} --- \cpp{compute\_ne} throws otherwise.  GCPM and IRI
describe only the inner-magnetospheric cold plasma; contributions beyond the
cutoff radius (default $4\,R_E$, see \cref{sec:iono-eikonal}) are neglected,
where the densities are orders of magnitude lower and the residual delay is well
below the noise floor of a lunar GNSS receiver.
\end{lupntnote}

\begin{lupntwarning}
GCPM cannot be evaluated for arbitrary future epochs: it requires historical
solar and geomagnetic indices, and the bundled \file{apf107.dat} and
\file{ig_rz.dat} tables end around 2024 with only a short projection.  Both
\cpp{iri\_sm} and \cpp{gcpm\_v24\_fortran} therefore step the requested year
backwards (up to 30 times) until the solar-index tables have coverage, so a
future epoch silently falls back to the most recent available activity rather
than returning NaN.  The thesis baseline sidesteps this by using
2025-01-01 12:00 UTC as an explicit proxy epoch for a 2027 scenario and by
sweeping $K_p \in \{1, 3, 5, 6.667, 9\}$ and
$R_{12} \in \{10, 50, 100, 150, 200\}$ \citep[Ch.~5]{iiyama2026thesis}.
\end{lupntwarning}

\section{Solar Indices: \texorpdfstring{$R_{12}$}{R12}, \texorpdfstring{$F_{10.7}$}{F10.7} and \texorpdfstring{$IG_{12}$}{IG12}}
\label{sec:iono-indices}

IRI is driven by three correlated solar indices.  When the user supplies a
sunspot number $R_{12}$ directly, \lupnt{} derives the other two from it with
the PHaRLAP relations \citep{cervera2014pharlap}, which are the same polynomial
fits used inside the IRI Fortran core (thesis Eqs.~5.19--5.20):
%
\begin{align}
  \label{eq:iono-f107}
  F_{10.7} &= 63.75 + 0.728\,R_{12} + \num{0.00089}\,R_{12}^2 , \\
  \label{eq:iono-ig12}
  IG_{12}  &= -12.349154 + 1.4683266\,R_{12} - \num{0.00267690893}\,R_{12}^2 .
\end{align}
%
$F_{10.7}$ (the \SI{10.7}{\centi\meter} solar radio flux, in solar flux units)
drives the topside and composition models; $IG_{12}$ is the
ionospheric-effective index used by IRI for the F2 critical frequency $f_oF_2$.
The derived triple is written into the IRI output array \code{oarr} at the
Fortran slots for $R_{12}$, $IG_{12}$ and $F_{10.7}$ before the model call, and
IRI2020 additionally receives $F_{10.7,81} = F_{10.7}$.

\begin{lstlisting}[language=cpp, caption={PHaRLAP solar-index relations and the
IRI user-input path.  Source:
\texttt{cpp/lupnt/environment/plasma/gcpm/iri\_interface.cc} ::
\code{iri\_sub}.}, label={lst:iono-indices}]
// cpp/lupnt/environment/plasma/gcpm/iri_interface.cc :: iri_sub  (IRI 2020 branch)
R12 = iri_option_2020.R12;
if (R12 > 0.0) {
  f107 = 63.75 + R12 * (0.728 + R12 * 0.00089);
  IG12 = -12.349154 + R12 * (1.4683266 - R12 * 2.67690893e-03);
  *(oarr + 32) = static_cast<float>(R12);
  *(oarr + 38) = static_cast<float>(IG12);
  *(oarr + 40) = static_cast<float>(f107);
  *(oarr + 45) = static_cast<float>(f107);   /* F10.7_81 is set to F10.7 */
}
iri_sub_2020_(iri_option_2020.jf2020, &jmag, &blatdf, &blongdf, &yyyy, &mddd,
              &dhourf, &aheight1f, &aheight2f, &delhf, outf, oarr);
\end{lstlisting}

The sign of \code{config.rz12} selects between user-supplied and
historical/projected indices following the IRI convention, and simultaneously
switches the foF2 storm model.  \cpp{trace\_ray} pushes \code{config.rz12} into
the active option object (\cpp{IRI2007Option} or \cpp{IRI2020Option}) and
re-derives the Fortran logical-flag array \code{jf} before tracing.

\begin{table}[htbp]
  \centering
  \caption{Interpretation of \code{RayTraceConfig::rz12} and the resulting IRI
  Fortran flags.  Flag indices are one-based, as in the IRI documentation.}
  \label{tab:iono-rz12}
  \small
  \begin{tabularx}{\textwidth}{lLl}
    \toprule
    \code{rz12} & Meaning & IRI flags \\
    \midrule
    $0 < R_{12} \le 200$ & User-supplied sunspot number; $F_{10.7}$ and $IG_{12}$ derived from \cref{eq:iono-f107,eq:iono-ig12} & \code{jf(17)}, \code{jf(25)}, \code{jf(27)}, \code{jf(32)} $=$ false; storm model off \\
    $-1$ & IRI historical / projected $R_{12}$, storm model on & all four $=$ true; \code{jf(26)} $=$ true \\
    $-2$ & IRI historical / projected $R_{12}$, storm model off (IRI2007 only) & \code{jf(26)} $=$ false \\
    other & Rejected --- \cpp{update\_jf\_2007} / \cpp{update\_jf\_2020} throws \code{"Invalid value for R12"} & --- \\
    \bottomrule
  \end{tabularx}
\end{table}

\section{GCPM Interpolation Surrogate}
\label{sec:iono-gcpm-surrogate}

GCPM is expensive because each call
performs $\sim$10 IRI evaluations.  For bulk work --- Monte-Carlo delay budgets,
data-assimilation experiments --- \lupnt{} provides an optional interpolation
surrogate (\code{RayTraceConfig::use\_gcpm\_surrogate}, default off;
\file{gcpm/gcpm\_surrogate.h}, \file{gcpm/gcpm\_surrogate.cc}).  GCPM's \emph{own}
output is pre-sampled on a seven-dimensional grid in its natural solar-magnetic
frame plus UT,
%
\begin{equation}
  \label{eq:iono-gcpm-surrogate}
  n_e = f\big(r,\ \phi_\mathrm{SM},\ \mathrm{MLT},\ \mathrm{UT},\ K_p,\ R_{12},\ \mathrm{doy}\big),
\end{equation}
%
stored as $\log n_e$ and interpolated multilinearly (with MLT/UT/season periodic),
so the model's plasmapause and region-blending logic are baked in.  It is used as a
\emph{hybrid}: above $r = 1.15\,R_E$, where GCPM is plasmasphere-organized and
interpolates cleanly, the table is queried; below it, where the sharp,
\emph{geographic} F-region is not representable on a solar-magnetic grid, the call
falls through to full GCPM.  The surrogate is
exact at grid nodes, gives a $\sim$2\% median error on ray-integrated lunar-link
delays, and is $14$--$200\times$ faster than full GCPM per call --- always more
than $10\times$.

\begin{lupntwarning}
The GCPM surrogate does \emph{not} reach sub-percent accuracy, and no practical
grid can: in solar-magnetic coordinates the ionosphere's geographic structure
(solar zenith, terminator, equatorial anomaly) becomes razor-sharp, and the
ionosphere--plasmasphere bridge keeps GCPM universal-time-dependent up to
$\sim$1.6\,$R_E$ at mid/high latitude.  A genuinely sub-percent fast GCPM would
require re-deriving a smooth parametric emulator ($N_{mF_2}$, $h_{mF_2}$, and scale
heights reconstructed into a Chapman profile).  Use the surrogate where a $\sim$2\% delay error
is acceptable in exchange for the speedup; leave it off for reference evaluations.
\end{lupntwarning}

\section{Geomagnetic Field: IGRF-14}
\label{sec:iono-igrf}

The magnetic field is required only by the second- and third-order delay terms
of \cref{sec:iono-refractive}.  \lupnt{} evaluates it with the IGRF-14
coefficient set \citep{alken2021igrf}, whose scalar potential is the truncated
spherical-harmonic series
%
\begin{equation}
  \label{eq:iono-igrf-potential}
  V(r,\vartheta,\lambda,t)
  = a \sum_{n=1}^{N} \sum_{m=0}^{n}
    \left(\frac{a}{r}\right)^{n+1}
    \Big[ g_n^m(t)\cos m\lambda + h_n^m(t)\sin m\lambda \Big]\,
    P_n^m(\cos\vartheta) ,
\end{equation}
%
where $a = \SI{6371.2}{\kilo\meter}$ is the IGRF reference radius,
$\vartheta$ is geocentric colatitude, $\lambda$ is east longitude, $P_n^m$ are
Schmidt quasi-normalized associated Legendre functions, $N = 13$ for the main
field, and the Gauss coefficients are linearly interpolated between five-year
epochs with a secular-variation extrapolation beyond the last epoch,
%
\begin{equation}
  g_n^m(t) = g_n^m(T_k) + (t - T_k)\,\dot{g}_n^m ,
\end{equation}
%
with $t$ the decimal year.  The local north/east/down (NEC) components follow
from $\vec{B} = -\nabla V$ as
%
\begin{equation}
  \label{eq:iono-bfield}
  X = \frac{1}{r}\frac{\partial V}{\partial \vartheta},
  \qquad
  Y = -\frac{1}{r\sin\vartheta}\frac{\partial V}{\partial \lambda},
  \qquad
  Z = \frac{\partial V}{\partial r},
  \qquad
  B = \sqrt{X^2 + Y^2 + Z^2} \ \ [\si{\nano\tesla}] .
\end{equation}
%
The projection of the field on the ray enters the delay through
%
\begin{equation}
  \label{eq:iono-costheta}
  \cos\theta = \frac{\vec{B}\cdot\uvec{s}}{B} ,
\end{equation}
%
i.e.\ the cosine of the angle between the local field line and the local
propagation direction.  The decimal year handed to the Fortran routine is
$t = \mathrm{year} + (\mathrm{doy}-1)/365.25$.

\begin{lstlisting}[language=cpp, caption={IGRF-14 evaluation and the ray/field
angle.  Sources: \texttt{cpp/lupnt/environment/plasma/igrf/igrf\_interface.cc} and
\texttt{cpp/lupnt/environment/plasma/tec/raytrace.cc} ::
\code{compute\_B}.}, label={lst:iono-igrf}]
// cpp/lupnt/environment/plasma/igrf/igrf_interface.cc :: igrf14
std::vector<double> igrf14(double lat_rad, double lon_rad, double r_km, double decimal_year) {
  int isv = 0;               // 0 = main field, not secular variation
  int itype = 2;             // 1 = geodetic, 2 = geocentric
  double x, y, z, f;
  double colat_deg = 90.0 - lat_rad * 180.0 / M_PI;
  double elong_deg = lon_rad * 180.0 / M_PI;
  if (elong_deg < 0) elong_deg += 360.0;
  igrf14syn_(&isv, &decimal_year, &itype, &r_km, &colat_deg, &elong_deg, &x, &y, &z, &f);
  return {x, y, z};          // north, east, vertical [nT]
}

// cpp/lupnt/environment/plasma/tec/raytrace.cc :: compute_B
cart_to_pol(pos_geo_arr, lat_rad, lon_rad, r);
double decimal_year = datetime.year + (datetime.doy - 1) / 365.25;
std::vector<double> B = igrf14(lat_rad, lon_rad, r, decimal_year);   // [nT]

// cpp/lupnt/environment/plasma/tec/raytrace.cc :: ray_derivative
Vec3d B_vec = compute_B(t_epoch, x, config, debug);   // in nT
double B_dot_s = B_vec.dot(s_hat);
B         = B_vec.norm();                             // in nT
cos_theta = B_dot_s / B;
\end{lstlisting}

\begin{lupntnote}
The delay coefficients hard-coded in \file{raytrace.cc} assume $B$ in
\textbf{nanotesla}, whereas the thesis writes the symbolic coefficients for $B$
in tesla.  This is the sole source of the numerical difference between the two
sets of constants; see \cref{tab:iono-coeffs}.
\end{lupntnote}

\begin{lupntwarning}
\cpp{igrf14} returns the field in the local north/east/vertical triad, but
\cpp{ray\_derivative} dots that triple directly with $\uvec{s}$, which is
expressed in geocentric Cartesian components.  The magnitude $B$ is unaffected
(it is frame-invariant), so the $n_e B^2(1+\cos^2\theta)$ part of the third-order
term is bounded correctly, but $\cos\theta$ as computed mixes the two triads and
is therefore only an order-of-magnitude proxy for the true field/ray angle.  The
terms affected are $d_{I2}$ and the $\cos^2\theta$ part of $d_{I3}$, which the
lunar characterization places several orders of magnitude below $d_{I1}$ and
which the GNSS default (\code{TEC\_ONLY}, \cref{sec:iono-gnss}) omits entirely.
Enable \code{compute\_higher\_order} for magnitude studies, not for
sub-centimetre error budgets.
\end{lupntwarning}

\section{Refractive Index, Group Delay, and TEC}
\label{sec:iono-refractive}

\subsection{Appleton--Hartree}
\label{sec:iono-appleton}

For a cold, magnetized, collisional plasma the complex phase refractive index of
a characteristic magnetoionic mode is given by the Appleton--Hartree relation
\citep{budden1985}
%
\begin{equation}
  \label{eq:iono-ah}
  n^2 = 1 -
  \frac{X}{\,1 - \mathrm{i}Z - \dfrac{Y_T^2}{2\,(1 - X - \mathrm{i}Z)}
        \pm \sqrt{\dfrac{Y_T^4}{4\,(1 - X - \mathrm{i}Z)^2} + Y_L^2}\,} ,
\end{equation}
%
with the dimensionless parameters
%
\begin{equation}
  \label{eq:iono-xyz}
  X = \frac{f_p^2}{f^2}, \qquad
  Y = \frac{f_g}{f}, \qquad
  Y_L = Y\cos\theta, \qquad
  Y_T = Y\sin\theta, \qquad
  Z = \frac{\nu}{2\pi f} ,
\end{equation}
%
where $\nu$ is the electron collision frequency.  The characteristic frequencies
are
%
\begin{equation}
  \label{eq:iono-fp}
  f_p^2 = \frac{n_e e^2}{4\pi^2 \epsilon_0 m_e} \approx 80.6\, n_e
  \ \ [\si{\square\hertz}] ,
  \qquad
  f_g = \frac{e B}{2\pi m_e} \approx 28\,B
  \ \ [\si{\hertz},\ B\ \text{in}\ \si{\nano\tesla}] ,
\end{equation}
%
with $e$ the elementary charge, $m_e$ the electron mass and $\epsilon_0$ the
vacuum permittivity.  The upper sign selects the ordinary mode and the lower
sign the extraordinary mode; a right-hand circularly polarized GNSS signal
propagating in the northern magnetic hemisphere follows the branch for which the
$f_g\cos\theta$ term appears with the sign written below.

At GNSS frequencies the ionosphere is weakly refracting and effectively
collisionless: $X \lesssim 10^{-4}$, $Y \lesssim 10^{-3}$ and $Z \approx 0$ above
the D region.  \lupnt{} therefore uses the quasi-longitudinal series expansion of
\cref{eq:iono-ah} truncated after the $f^{-4}$ term (thesis Eqs.~5.1--5.2):
%
\begin{align}
  \label{eq:iono-phase-index}
  n &= 1
     - \frac{f_p^2}{2f^2}
     - \frac{f_p^2 f_g \cos\theta}{2f^3}
     - \frac{f_p^2}{4f^4}\!\left[\frac{f_p^2}{2} + f_g^2\left(1+\cos^2\theta\right)\right] , \\[4pt]
  \label{eq:iono-group-index}
  n_\mathrm{gr} &= 1
     + \frac{f_p^2}{2f^2}
     + \frac{f_p^2 f_g \cos\theta}{f^3}
     + \frac{3f_p^2}{4f^4}\!\left[\frac{f_p^2}{2} + f_g^2\left(1+\cos^2\theta\right)\right] .
\end{align}
%
The group index follows from the phase index through
$n_\mathrm{gr} = \dd(fn)/\dd f$, which is what produces the factors
$1 \to 2 \to 3$ on the successive orders.  The delay is dispersive, scaling as
$f^{-2}$ at leading order; this is what makes the dual-frequency
ionosphere-free combination of \cref{chap:gnss-measurements} possible.

\subsection{Group delay along the path}
\label{sec:iono-groupdelay}

Integrating $n_\mathrm{gr}-1$ along the ray gives the excess group path (thesis
Eqs.~5.4--5.7):
%
\begin{equation}
  \label{eq:iono-delay-orders}
  d_{I_\mathrm{gr}} = \int_\Gamma \left(n_\mathrm{gr}-1\right)\dd s
    = \frac{p}{f^2} + \frac{q}{f^3} + \frac{u}{f^4} ,
\end{equation}
%
with, in SI units and $B$ in tesla,
%
\begin{align}
  \label{eq:iono-p}
  p &= 40.3 \int_\Gamma n_e\,\dd s
     = 40.3\,\mathrm{TEC}
     = 40.3\left(\mathrm{TEC}_\mathrm{LOS} + \Delta\mathrm{TEC}_\mathrm{bend}\right) , \\
  \label{eq:iono-q}
  q &= -\num{2.2566e12} \int_\Gamma n_e\, B \cos\theta\,\dd s , \\
  \label{eq:iono-u}
  u &= 2437 \int_\Gamma n_e^2\,\dd s
     + \num{4.74e22} \int_\Gamma n_e\, B^2 \left(1+\cos^2\theta\right)\dd s .
\end{align}
%
The first-order coefficient is $40.3 = \tfrac{1}{2}\cdot 80.6$ and carries units
of \si{\meter\cubed\per\square\second}; the second-order coefficient is
$\num{2.2566e12} = 7527\,c$, the familiar higher-order ionospheric constant of
the GNSS literature.  The integral of $n_e$ along the path is the total electron
content,
%
\begin{equation}
  \label{eq:iono-tecu}
  \mathrm{TEC} = \int_\Gamma n_e\,\dd s\ \ [\si{\per\square\meter}],
  \qquad
  \SI{1}{\TECU} = \num{1e16}\,\si{\per\square\meter},
  \qquad
  \SI{1}{\TECU} \ \Longrightarrow\ \SI{0.162}{\meter}\ \text{at L1} .
\end{equation}
%
Because $\Gamma$ is curved, the TEC splits into a line-of-sight part
$\mathrm{TEC}_\mathrm{LOS}$, obtained by sampling the same density field along
the straight chord, and a bending increment
$\Delta\mathrm{TEC}_\mathrm{bend} = \mathrm{TEC} - \mathrm{TEC}_\mathrm{LOS}$.
For terrestrial and LEO users the second- and third-order terms remain below
about \SI{1}{\percent} of the first order; for the grazing lunar geometry the
characterization of \citep[Tab.~5.2]{iiyama2026thesis} finds them several orders
of magnitude below $d_{I1}$.

\subsection{Sign conventions}
\label{sec:iono-signs}

Three sign conventions must be kept apart.

\begin{enumerate}[nosep]
  \item \textbf{Phase versus group.}  $n < 1$ and $n_\mathrm{gr} > 1$: the code
        observable is \emph{delayed} and the carrier phase is \emph{advanced} by
        the same first-order amount.  \cref{eq:iono-delay-orders} is the group
        (code) delay and is positive.
  \item \textbf{The stored $q$.}  The code accumulates $q$ with a
        \emph{positive} coefficient \code{q\_coeff}.  The physical sign is
        carried by $\cos\theta$ (which changes sign when the ray reverses
        relative to the field) and, in the phase index, by the explicit minus in
        $n = 1 - (p/f^2 + q/2f^3 + u/3f^4)$.  The thesis instead absorbs the sign
        into the symbolic definition of $q$ in Eq.~(5.6), hence the leading minus
        in \cref{eq:iono-q}.
  \item \textbf{Phase versus group factors.}  The phase index uses
        $q/(2f^3)$ and $u/(3f^4)$; the group-delay accumulators use $q/f^3$ and
        $u/f^4$.  Both are consistent with
        \cref{eq:iono-phase-index,eq:iono-group-index} --- the group form simply
        omits the $\tfrac{1}{2}$ and $\tfrac{1}{3}$ phase factors.
\end{enumerate}

\begin{lstlisting}[language=cpp, caption={Phase index used to bend the ray, and
the three group-delay accumulators.  Source:
\texttt{cpp/lupnt/environment/plasma/tec/raytrace.cc} ::
\code{refractive\_index\_neB}, \code{ray\_derivative}.}, label={lst:iono-index}]
// cpp/lupnt/environment/plasma/tec/raytrace.cc
static const double p_coeff  = 40.3;     // first-order
static const double q_coeff  = 2.2566e3; // second-order  (for B in nT)
static const double u1_coeff = 2437;     // third-order, ne^2 part
static const double u2_coeff = 4.74e4;   // third-order, B^2 part (for B in nT)

// :: refractive_index_neB  -- the phase index that bends the ray
double p = p_coeff * ne_m3;
double q = 0.0, u = 0.0;
if (compute_q) {
  q = q_coeff * ne_m3 * B * cos_theta;
  u = u1_coeff * ne_m3 * ne_m3 + u2_coeff * ne_m3 * B * B * (1 + cos_theta * cos_theta);
}
double n = 1 - (p / f2 + q / (2 * f3) + u / (3 * f4));   // phase refractive index

// :: ray_derivative -- per-arc group-delay accumulators (s is in km, hence *1000)
dz(7) = ne_m3 * 1000;                                  // d(TEC)/ds       [m^-2 per km]
dz(8) = q_coeff * ne_m3 * B * cos_theta * 1000;        // second-order integrand
dz(9) = 1000 * (u1_coeff * ne_m3 * ne_m3
              + u2_coeff * ne_m3 * B * B * (1 + cos_theta * cos_theta));  // third order
\end{lstlisting}

\begin{table}[htbp]
  \centering
  \caption{Delay coefficients: code constants (with $B$ in \si{\nano\tesla})
  versus the SI symbolic values of \citep[Ch.~5]{iiyama2026thesis} (with $B$ in
  \si{\tesla}).}
  \label{tab:iono-coeffs}
  \small
  \begin{tabular}{llll}
    \toprule
    Term & Code constant & Value ($B$ in \si{\nano\tesla}) & SI value ($B$ in \si{\tesla}) \\
    \midrule
    $p$, first order      & \code{p\_coeff}  & $40.3$          & $40.3$ \\
    $q$, second order     & \code{q\_coeff}  & $\num{2.2566e3}$ & $\num{2.2566e12}$ \quad ($\times 10^{-9}$) \\
    $u$, third order, $n_e^2$ & \code{u1\_coeff} & $2437$      & $2437$ \\
    $u$, third order, $B^2$   & \code{u2\_coeff} & $\num{4.74e4}$ & $\num{4.74e22}$ \quad ($\times 10^{-18}$) \\
    \bottomrule
  \end{tabular}
\end{table}

\section{Vector Eikonal Ray Tracing}
\label{sec:iono-eikonal}

\subsection{Governing equations and state}
\label{sec:iono-state}

In the geometric-optics limit the signal path obeys the vector eikonal (ray)
equations, integrated with respect to arc length $s$ (thesis Eqs.~5.11--5.13):
%
\begin{equation}
  \label{eq:iono-eikonal}
  \frac{\dd t}{\dd s} = \frac{n}{c},
  \qquad
  \frac{\dd \vec{r}}{\dd s} = \uvec{s},
  \qquad
  \frac{\dd \uvec{s}}{\dd s}
    = \frac{\nabla n - \uvec{s}\,(\uvec{s}\cdot\nabla n)}{n} ,
\end{equation}
%
subject to the two-point boundary conditions
$\vec{r}(t_\mathrm{tx}) = \vec{r}_\mathrm{tx}$ and
$\vec{r}(t_\mathrm{rx}) = \vec{r}_\mathrm{rx}$.  The right-hand side of the third
equation is the component of $\nabla n$ perpendicular to the ray, normalized by
$n$: the ray bends towards increasing refractive index, and the parallel
component --- which would only rescale $\lVert\uvec{s}\rVert$ --- is projected
out, keeping $\uvec{s}$ on the unit sphere.

The integrated state has ten components,
%
\begin{equation}
  \label{eq:iono-statevec}
  \vec{z} = \big[\,\underbrace{\vec{r}}_{0..2},\ \underbrace{\uvec{s}}_{3..5},\
                  \underbrace{\Delta t}_{6},\ \underbrace{\mathrm{TEC}}_{7},\
                  \underbrace{I_2}_{8},\ \underbrace{I_3}_{9}\,\big]
       \in \mathbb{R}^{10},
\end{equation}
%
where $\Delta t = t - t_\mathrm{tx}$ and $I_2$, $I_3$ are the second- and
third-order delay integrals of \cref{eq:iono-q,eq:iono-u} (without the $f^{-3}$,
$f^{-4}$ factors, which are applied once at the end).  The refractive-index
gradient is a central finite difference with step \code{config.gradn\_dx}
(default \SI{1}{\kilo\meter}):
%
\begin{equation}
  \label{eq:iono-gradn}
  \left[\nabla n\right]_i
  = \frac{n(\vec{x} + \delta\,\uvec{e}_i) - n(\vec{x} - \delta\,\uvec{e}_i)}{2\delta},
  \qquad i = 1,2,3, \qquad \delta = \code{gradn\_dx} .
\end{equation}
%
Each gradient evaluation therefore costs six GCPM calls, which dominates the
runtime and motivates both the stored-gradient reuse of
\cref{sec:iono-correction} and the batch precompute path of
\cref{sec:iono-precompute}.

\begin{lstlisting}[language=cpp, caption={Eikonal right-hand side.  Source:
\texttt{cpp/lupnt/environment/plasma/tec/raytrace.cc} ::
\code{ray\_derivative}.}, label={lst:iono-derivative}]
// cpp/lupnt/environment/plasma/tec/raytrace.cc :: ray_derivative
Vec3d dshat_ds = Vec3d::Zero();
if (config.straight_ray) dshat_ds.setZero();
else                     dshat_ds = (grad_n - s_hat * grad_n.dot(s_hat)) / n;   // Eq. (5.13)

VecXd dz(STATE_SIZE);
dz.segment<3>(0) = s_hat;      // dr/ds   = s_hat
dz.segment<3>(3) = dshat_ds;   // ds_hat/ds
// We only track the time increase due to distance here due to machine precision
// in tracking TEC delays (= p_coeff * ne_m3 / (freq_Hz * freq_Hz) / C = ne_m3 * 1e-23)
// -> We do not want to keep epochs and TEC delays in the same vector
dz(6) = 1.0 / C;               // dt/ds  ~ 1/c (vacuum; the plasma delay is added later)
\end{lstlisting}

\begin{lupntnote}
The geometric propagation time is accumulated as $\dd t/\dd s = 1/c$ (vacuum)
rather than $n/c$, and the plasma contribution is added once at the end as
$(d_{I1} + d_{I2} + d_{I3})/c$.  This is a deliberate floating-point choice: the
per-step TEC time increment is of order $n_e \cdot 10^{-23}\,\si{\second}$, far
below the resolution of a double holding an epoch of order $10^{8}\,\si{\second}$
from J2000, so co-accumulating the two would lose the delay entirely.  The result
is algebraically identical to integrating $\dd t/\dd s = n/c$.
\end{lupntnote}

\subsection{Integration and the adaptive step schedule}
\label{sec:iono-integration}

\code{config.integ\_method} selects between the explicit Euler step
%
\begin{equation}
  \label{eq:iono-euler}
  \vec{z}_{k+1} = \vec{z}_k + h\,\vec{F}(s_k, \vec{z}_k),
  \qquad s_{k+1} = s_k + h ,
\end{equation}
%
and the classical fourth-order Runge--Kutta step \citep{hairer1993}
%
\begin{equation}
  \label{eq:iono-rk4}
  \begin{gathered}
    \vec{k}_1 = \vec{F}(s_k, \vec{z}_k), \qquad
    \vec{k}_2 = \vec{F}\!\left(s_k + \tfrac{h}{2},\ \vec{z}_k + \tfrac{h}{2}\vec{k}_1\right), \\
    \vec{k}_3 = \vec{F}\!\left(s_k + \tfrac{h}{2},\ \vec{z}_k + \tfrac{h}{2}\vec{k}_2\right), \qquad
    \vec{k}_4 = \vec{F}\!\left(s_k + h,\ \vec{z}_k + h\,\vec{k}_3\right), \\
    \vec{z}_{k+1} = \vec{z}_k
      + \frac{h}{6}\left(\vec{k}_1 + 2\vec{k}_2 + 2\vec{k}_3 + \vec{k}_4\right) .
  \end{gathered}
\end{equation}
%
RK4 costs four density/gradient evaluations per step (i.e.\ 28 GCPM calls with
the finite-difference gradient) and correspondingly quadruples the storage
width; Euler is the default because the adaptive schedule below already keeps
the local truncation error well inside the model error of GCPM itself.  After
each step $\uvec{s}$ is renormalized, which removes the drift that the
projection in \cref{eq:iono-eikonal} leaves at $O(h^2)$.

The step size is not error-controlled but keyed on altitude
$\mathrm{alt} = r - R_E$, which is an excellent proxy for the density scale
height:
%
\begin{equation}
  \label{eq:iono-step}
  h(\mathrm{alt}) =
  \begin{cases}
    \SI{10}{\kilo\meter},  & \mathrm{alt} < \SI{1000}{\kilo\meter}, \\
    \SI{20}{\kilo\meter},  & \SI{1000}{\kilo\meter} \le \mathrm{alt} < \SI{4000}{\kilo\meter}, \\
    \SI{100}{\kilo\meter}, & \mathrm{alt} \ge \SI{4000}{\kilo\meter} .
  \end{cases}
\end{equation}
%
The schedule is enabled by \code{config.use\_adaptive\_step} (default true) and
overrides \code{config.step\_size} within the plasmasphere; the final partial
step is clipped so that $s$ lands exactly on $s_f$.  Because the schedule can
reduce the step to \SI{10}{\kilo\meter} even when \code{config.step\_size} is
larger, the per-step storage matrix is allocated for the worst case
$\lceil s_f/\SI{10}{\kilo\meter}\rceil + 2$ rows and trimmed afterwards.

\begin{lstlisting}[language=cpp, caption={Altitude-keyed step schedule.  Source:
\texttt{cpp/lupnt/environment/plasma/tec/raytrace.cc} ::
\code{adjust\_stepsize}.}, label={lst:iono-step}]
// cpp/lupnt/environment/plasma/tec/raytrace.cc :: adjust_stepsize
double adjust_stepsize(double r, const VecXd& store_vec) {
  double step_size = 20.0;   // Default step size in km
  double alt = r - RE;       // Altitude in km
  if (alt < 1000)      step_size = 10.0;    // Smaller step size for closer distances
  else if (alt < 4000) step_size = 20.0;    // Default step size for medium distances
  else                 step_size = 100.0;   // Maximum step size for very far distances
  return step_size;
}
\end{lstlisting}

\subsection{Cutoff split and the terminal shot}
\label{sec:iono-cutoff}

The path is split into a plasmasphere section and a vacuum section at the cutoff
radius \code{config.cutoff\_r} (default
$4\,R_E \approx \SI{25484}{\kilo\meter}$).  Beyond the cutoff the density is
taken as zero, so $\nabla n = \vec{0}$ and the ray is exactly straight; there is
no point integrating it.  When the marching state first satisfies
$r \ge r_\mathrm{cut}$ \emph{and} $r > r_\mathrm{prev}$ (the second test prevents
an inbound ray from terminating on entry), the tracer projects the remaining
displacement onto the exit direction (thesis Eqs.~5.17--5.18):
%
\begin{equation}
  \label{eq:iono-shoot}
  \Delta s = \left(\vec{x}_\mathrm{rx} - \vec{x}_\mathrm{exit}\right)\cdot\uvec{s}_\mathrm{exit},
  \qquad
  s_f = s_\mathrm{exit} + \Delta s,
  \qquad
  \vec{x}_f = \vec{x}_\mathrm{exit} + \Delta s\,\uvec{s}_\mathrm{exit} .
\end{equation}
%
$\Delta s$ is the least-squares closest approach along the exit ray, so the
residual $\vec{x}_f - \vec{x}_\mathrm{rx}$ is exactly the component of the miss
perpendicular to $\uvec{s}_\mathrm{exit}$ --- which is precisely the quantity the
direction correction of \cref{sec:iono-correction} drives to zero.  The
propagation time is extended by the vacuum flight time $\Delta s / c$.

\begin{lstlisting}[language=cpp, caption={Cutoff exit and terminal shot.  Source:
\texttt{cpp/lupnt/environment/plasma/tec/raytrace.cc} ::
\code{propagate\_ray}.}, label={lst:iono-shoot}]
// cpp/lupnt/environment/plasma/tec/raytrace.cc :: propagate_ray
if ((r >= cutoff_r) && (r > r_prev)) {
  // Find sf that minimizes |fpos - x2|, where fpos = pos_now + dir_now * (sf - s)
  Vec3d delta_x2 = x2 - pos_now;
  double delta_s = delta_x2.dot(dir_now);
  sf = s + delta_s;
  Vec3d final_pos = pos_now + dir_now * (sf - s);
  double computed_prop_time = z(6) + (sf - s) / C;
  z.segment<3>(0) = final_pos;
  z(6) = computed_prop_time;
  cutoff_hit_p = true;
  break;
}
\end{lstlisting}

\subsection{Receiver-side convergence iteration}
\label{sec:iono-correction}

\Cref{eq:iono-eikonal} is an initial-value problem, but the physical problem is a
two-point boundary-value problem: the unknown launch direction $\uvec{s}_0$ must
be chosen so that the ray lands on the receiver.  \lupnt{} solves it by shooting.
Parameterizing the launch direction by azimuth and elevation,
$\uvec{s}_0 = \uvec{s}(\alpha, \epsilon)$, the objective is the terminal miss
distance
%
\begin{equation}
  \label{eq:iono-miss}
  J(\alpha,\epsilon)
  = \left\lVert \bar{\vec{x}}_\mathrm{rx}(\alpha,\epsilon) - \vec{x}_\mathrm{rx} \right\rVert ,
\end{equation}
%
where $\bar{\vec{x}}_\mathrm{rx}$ is the terminal position produced by
\cpp{propagate\_ray}.  The iteration is initialized with the straight-line
direction and chord length,
$\uvec{s}_0 = (\vec{x}_2 - \vec{x}_1)/\rho$ and $s_f = \rho$, and alternates two
updates for at most 20 coarse iterations:

\begin{enumerate}[nosep]
  \item \textbf{Propagation-time refresh} (\cpp{correct\_proptime}).  The
        propagation time is reset to include the plasma delay accumulated on the
        current path,
        %
        \begin{equation}
          \label{eq:iono-proptime}
          \tau \leftarrow \Delta t
            + \frac{1}{c}\left(\frac{p}{f^2} + \frac{q}{f^3} + \frac{u}{f^4}\right) ,
        \end{equation}
        %
        which in turn moves $t_\mathrm{tx} = t_\mathrm{rx} - \tau$ and hence the
        epoch at which GCPM is sampled.  This is the light-time iteration of
        thesis Eq.~(5.16), consistent with \cref{chap:gnss-measurements}.
  \item \textbf{Direction re-solve} (\cpp{correct\_dir\_neldermead}).  $J$ is
        minimized over $(\alpha, \epsilon)$ with a two-dimensional Nelder--Mead
        simplex.
\end{enumerate}

The coarse loop terminates as soon as one of
%
\begin{equation}
  \left|J_k - J_{k-1}\right| < \SI{1}{\meter},
  \qquad
  J_k > J_{k-1},
  \qquad
  J_k \le \code{corr\_tol}
\end{equation}
%
holds; if the error increased, the previous direction, state and stored
derivative matrix are restored.  To keep the cost tractable, the coarse pass
runs with \code{use\_precomputed\_vals = true}: it replays the per-step
refractive index, gradient, $n_e$, $B$ and $\cos\theta$ stored in
\code{store\_mat} on the previous propagation, together with the step sizes that
produced them, instead of re-querying GCPM at every point.  Only the geometry is
re-integrated.  The stored rows are trimmed to those actually written, so a
perturbed ray that runs longer than the reference simply stops at the last stored
row.

An optional finer pass (\code{config.fine\_correction}) re-solves with GCPM
evaluated at every step, either with Nelder--Mead again or with a
finite-difference Newton update:
%
\begin{equation}
  \label{eq:iono-newton}
  \mat{J} = \frac{\partial \bar{\vec{x}}_\mathrm{rx}}{\partial (\alpha,\epsilon)}
          \in \mathbb{R}^{3\times 2},
  \qquad
  \begin{bmatrix}\Delta\alpha \\ \Delta\epsilon\end{bmatrix}
    = \mat{J}^{+}\left(\vec{x}_\mathrm{rx} - \bar{\vec{x}}_\mathrm{rx}\right),
\end{equation}
%
where $\mat{J}$ is built by central differences with a perturbation of
$10^{-7}\,\si{\radian}$ and $\mat{J}^{+}$ is the thin-SVD pseudo-inverse (the
system is over-determined because the terminal shot of \cref{eq:iono-shoot}
already removed the along-track component).  The step is applied with unit
relaxation.  \Cref{tab:iono-neldermead} lists the simplex settings.

\begin{table}[htbp]
  \centering
  \caption{Nelder--Mead settings in \cpp{correct\_dir\_neldermead}.  $J_0$ is the
  terminal miss of the incoming solution, in \si{\kilo\meter}.}
  \label{tab:iono-neldermead}
  \small
  \begin{tabularx}{\textwidth}{lLL}
    \toprule
    Setting & Coarse pass (\code{coarse\_mode = true}) & Fine pass (\code{coarse\_mode = false}) \\
    \midrule
    Density/gradient source & replayed from \code{store\_mat} & re-evaluated from GCPM/IGRF \\
    Initial simplex edge     & $10^{-4}\,\si{\radian}$ & $10^{-6}$, $10^{-7}$ or $10^{-8}\,\si{\radian}$ for $J_0 > 0.1$, $J_0 \le 0.1$, $J_0 \le 0.01$ \\
    Simplex tolerance        & $10^{-8}$ & $10^{-4}$ \\
    Objective tolerance      & $10^{-6}\,\si{\kilo\meter}$ & \code{corr\_tol}$/1000$ \si{\kilo\meter} \\
    Iteration cap            & 200 & 50 \\
    \bottomrule
  \end{tabularx}
\end{table}

\begin{lstlisting}[language=cpp, caption={Coarse convergence loop: propagate,
refresh the propagation time, re-solve the launch direction.  Source:
\texttt{cpp/lupnt/environment/plasma/tec/raytrace.cc} ::
\code{trace\_ray}.}, label={lst:iono-traceray}]
// cpp/lupnt/environment/plasma/tec/raytrace.cc :: trace_ray
Vec3d dir = (x2 - x1).normalized();
double L  = (x2 - x1).norm();
double prop_time = L / C;     // straight-line light time as the initial guess
double sf = L;

int n_iter = 20;              // Number of iterations for coarse correction
for (int iter = 0; iter < n_iter; ++iter) {
  zf = propagate_ray(prop_time, dir, sf, x1, x2, epoch_utc_rx, config, store_mat, false);
  pos_err = (zf.segment<3>(0) - x2).norm();                      // [km]

  if ((abs(pos_err - pos_err_prev) < 1e-3) || (pos_err > pos_err_prev)
      || (pos_err <= config.corr_tol / 1000)) {
    if (pos_err > pos_err_prev) { dir = dir_prev; zf = zf_prev; store_mat = store_mat_prev; }
    break;
  }
  correct_proptime(prop_time, dir, sf, x1, x2, epoch_utc_rx, config, zf);
  pos_err_prev = pos_err; dir_prev = dir; zf_prev = zf; store_mat_prev = store_mat;
  correct_dir_neldermead(prop_time, dir, sf, x1, x2, epoch_utc_rx, config, zf, store_mat,
                         /*coarse_mode=*/true, debug_corr);
}
\end{lstlisting}

Without the correction the terminal miss is \SIrange{10}{100}{\kilo\meter} for a
grazing lunar link, purely because of bending; with it the miss is driven to the
sub-metre to \SI{100}{\meter} range, depending on whether the fine pass is
enabled.  The residual is reported in \code{PathProfile::corr\_final\_pos\_err}
and \code{PathProfile::corr\_final\_time\_err} so that a caller can reject
non-converged rays.

\begin{lupntwarning}
\code{config.correction\_method} must be \code{"neldermead"} or \code{"newton"};
\cpp{trace\_ray} throws otherwise.  The struct default is \code{"grid"}, which
\emph{would} throw, so any caller that enables \code{correction} must set this
field explicitly.  The thrown message text
(\code{"Supported methods are 'naive' and 'newton'"}) is stale with respect to
the check it accompanies.
\end{lupntwarning}

\section{Delay Decomposition}
\label{sec:iono-decomposition}

Once the corrected path has been integrated, \cpp{compute\_path\_profile}
converts the accumulators $z_7$, $z_8$, $z_9$ and the geometry into metres.  The
excess geometric path caused by bending is
%
\begin{equation}
  \label{eq:iono-dlen}
  d_\mathrm{len} = \int_\Gamma \dd s - \rho = s_f - \rho,
  \qquad
  \rho = \left\lVert \vec{r}_\mathrm{tx} - \vec{r}_\mathrm{rx} \right\rVert ,
\end{equation}
%
which is non-negative because the straight chord is the shortest path between the
endpoints.  The total group delay relative to the chord is then the sum of five
physically distinct terms (thesis Eqs.~5.8--5.10):
%
\begin{equation}
  \label{eq:iono-dtotal}
  d_\mathrm{total}
    = d_{I_\mathrm{gr}} + d_\mathrm{len}
    = \underbrace{40.3\,\mathrm{TEC}_\mathrm{LOS}}_{d_{I1,\mathrm{LOS}}}
    + \underbrace{q/f^3}_{d_{I2}}
    + \underbrace{u/f^4}_{d_{I3}}
    + \underbrace{40.3\,\Delta\mathrm{TEC}_\mathrm{bend}}_{d_{I1,B}}
    + d_\mathrm{len} .
\end{equation}
%
$d_{I1,\mathrm{LOS}}$ is what a straight-line TEC model would predict;
$d_{I1,B}$ is the additional dispersive delay incurred because the bent path
samples a different (generally denser, lower-altitude) part of the ionosphere;
$d_\mathrm{len}$ is the non-dispersive geometric lengthening.  The straight-line
reference $\mathrm{TEC}_\mathrm{LOS}$ is computed by
\cpp{compute\_tec\_straight}, which marches the same density field along the
chord with fixed steps of \code{config.step\_size} and zeroes any sample outside
the cutoff radius.

\begin{lstlisting}[language=cpp, caption={Assembly of the delay terms.  Source:
\texttt{cpp/lupnt/environment/plasma/tec/raytrace.cc} ::
\code{compute\_path\_profile}.}, label={lst:iono-profile}]
// cpp/lupnt/environment/plasma/tec/raytrace.cc :: compute_path_profile
double tec_straight = z(7);
if (config.straight_ray == false)
  tec_straight = compute_tec_straight(t_tx, initial_pos, final_pos, config);

double dist                  = (final_pos - initial_pos).norm();       // chord [km]
double tec_delay_m           = p_coeff * z(7) / freq2;                 // 40.3 * TEC_bent / f^2
double tec_delay_straight_m  = p_coeff * tec_straight / freq2;
double second_delay_m        = z(8) / freq3;
double third_delay_m         = z(9) / freq4;

double C_ms = C * 1000;                                                // [m/s]
double prop_time_total = z(6) + (tec_delay_m + second_delay_m + third_delay_m) / C_ms;

pp.dist_bend_m      = (sf - dist) * 1000;                              // excess path length [m]
pp.tec_delay_m      = tec_delay_m;
pp.tec_delay_bend_m = tec_delay_m - tec_delay_straight_m;
pp.total_delay_m    = pp.dist_bend_m + pp.tec_delay_m
                    + pp.second_delay_m + pp.third_delay_m;
pp.tecu             = z(7) / TECU;
pp.tecu_bend        = (z(7) - tec_straight) / TECU;
\end{lstlisting}

\begin{table}[htbp]
  \centering
  \caption{\cpp{PathProfile} scalar outputs and their mapping onto
  \cref{eq:iono-dtotal}.}
  \label{tab:iono-pathprofile}
  \small
  \begin{tabularx}{\textwidth}{llL}
    \toprule
    Field & Symbol & Meaning \\
    \midrule
    \code{tec\_delay\_m}      & $d_{I1,\mathrm{LOS}} + d_{I1,B}$ & First-order delay $40.3\,\mathrm{TEC}/f^2$ along the \emph{bent} path [\si{\meter}] \\
    \code{tec\_delay\_bend\_m}& $d_{I1,B}$        & Bending increment $40.3\,\Delta\mathrm{TEC}_\mathrm{bend}/f^2$ [\si{\meter}] \\
    \code{second\_delay\_m}   & $d_{I2}$          & $q/f^3$ [\si{\meter}] \\
    \code{third\_delay\_m}    & $d_{I3}$          & $u/f^4$ [\si{\meter}] \\
    \code{dist\_bend\_m}      & $d_\mathrm{len}$  & Excess geometric path $(s_f - \rho)$ [\si{\meter}] \\
    \code{total\_delay\_m}    & $d_\mathrm{total}$& Sum of the four rows above [\si{\meter}] \\
    \code{tecu}               & $\mathrm{TEC}$    & Bent-path TEC [\si{\TECU}] \\
    \code{tecu\_bend}         & $\Delta\mathrm{TEC}_\mathrm{bend}$ & Bent minus straight TEC [\si{\TECU}] \\
    \code{sf}                 & $s_f$             & Total traced arc length [\si{\kilo\meter}] \\
    \code{dist\_straight\_km} & $\rho$            & Chord length [\si{\kilo\meter}] \\
    \code{max\_sep\_line\_m}  & ---               & Maximum perpendicular separation of $\Gamma$ from the chord [\si{\meter}] \\
    \code{prop\_time\_total}  & $\tau$            & Total propagation time including plasma delay [\si{\second}] \\
    \code{corr\_final\_pos\_err} & ---            & Residual terminal miss $\bar{\vec{x}}_\mathrm{rx}-\vec{x}_\mathrm{rx}$ [\si{\kilo\meter}] \\
    \code{corr\_final\_time\_err}& ---            & Residual propagation-time inconsistency [\si{\second}] \\
    \bottomrule
  \end{tabularx}
\end{table}

The per-step arrays \code{s}, \code{r}, \code{pos\_eci}, \code{tec\_section},
\code{az\_dir}, \code{el\_dir} and \code{dist\_to\_line} carry the sampled
profile itself and are what the characterization plots of
\citep[Ch.~5]{iiyama2026thesis} are built from.

\section{Configuration Knobs}
\label{sec:iono-config}

\begin{table}[htbp]
  \centering
  \caption{\cpp{RayTraceConfig} (\texttt{cpp/lupnt/environment/plasma/tec/raytrace.h}).}
  \label{tab:iono-config}
  \small
  \begin{tabularx}{\textwidth}{llL}
    \toprule
    Field & Default & Effect \\
    \midrule
    \code{freq\_Hz}   & NaN & Carrier frequency [\si{\hertz}].  The constants \code{freq\_L1} $=\SI{1575.42}{\mega\hertz}$, \code{freq\_L2} $=\SI{1227.60}{\mega\hertz}$, \code{freq\_L5} $=\SI{1176.45}{\mega\hertz}$ are provided.  All delays scale as $f^{-2}$ at leading order, so L5 sees $\approx 1.8\times$ the L1 delay.  Must be set. \\
    \code{step\_size} & $10.0$ & Nominal integrator step [\si{\kilo\meter}]; also the fixed step of the straight-line TEC reference. \\
    \code{use\_adaptive\_step} & true & Enable the 10/20/100 \si{\kilo\meter} altitude schedule of \cref{eq:iono-step}. \\
    \code{integ\_method} & \code{"Euler"} & \code{"Euler"} or \code{"RK4"} (case-insensitive for RK4); anything else throws. \\
    \code{correction} & true & Run the bent-path launch-direction correction of \cref{sec:iono-correction}. \\
    \code{fine\_correction} & false & Run the additional fine pass with GCPM re-evaluated at every step. \\
    \code{correction\_method} & \code{"grid"} & Must be set to \code{"neldermead"} or \code{"newton"} when \code{correction} is on. \\
    \code{corr\_tol} & $1.0$ & Terminal miss tolerance [\si{\meter}]. \\
    \code{cutoff\_r} & $4R_E$ & Plasmasphere/vacuum boundary [\si{\kilo\meter}]. \\
    \code{gradn\_dx} & $1.0$ & Central-difference step for $\nabla n$ [\si{\kilo\meter}]. \\
    \code{kp} & $-1$ & Geomagnetic $K_p$; must be set $\ge 0$ or \cpp{compute\_ne} throws. \\
    \code{rz12} & $-1.0$ & Solar $R_{12}$; see \cref{tab:iono-rz12}. \\
    \code{compute\_higher\_order} & true & Include $d_{I2}$, $d_{I3}$ (requires the IGRF evaluation). \\
    \code{straight\_ray} & false & Force $\nabla n = \vec{0}$ and skip the coarse correction, producing the straight-line TEC baseline. \\
    \code{use\_fortran\_gcpm} & true & Use the Fortran GCPM v2.4 backend instead of the C++ reimplementation (GCPM backend only). \\
    \bottomrule
  \end{tabularx}
\end{table}

The electron-density backend itself is not a \cpp{RayTraceConfig} field but a
process-global selected with \cpp{set\_iono\_model},
in the same way \cpp{set\_iri\_model} selects the IRI submodel.

\section{Coupling into the GNSS Measurement Model}
\label{sec:iono-gnss}

When \code{options.apply\_ionosphere\_plasma\_delay} is true and ray-trace
options have been registered through
\cpp{SetIonospherePlasmaRayTraceOptions}, the GNSS model calls the tracer once
per link.

\begin{lstlisting}[language=cpp, caption={Per-link tracer invocation.  Source:
\texttt{cpp/lupnt/measurements/gnss\_measurement.cc} ::
\code{ComputeIonospherePlasmaRayTraceDelay}.}, label={lst:iono-gnss}]
// cpp/lupnt/measurements/gnss_measurement.cc :: ComputeIonospherePlasmaRayTraceDelay
pecsim::RayTraceConfig config = ray_opts.config;
if (ray_opts.use_channel_frequency) config.freq_Hz = GnssFrequencyHz(freq);
if (ray_opts.iri_model != pecsim::IRIModel::NONE) pecsim::set_iri_model(ray_opts.iri_model);

Real t_raytrace = ConvertGnssEpoch(receive_time, options_.receive_time_scale,
                                   ray_opts.raytrace_epoch_scale);
Real t_tdb      = ConvertGnssEpoch(receive_time, options_.receive_time_scale, Time::TDB);

Vec3 r_rx_ray = r_rx, r_tx_ray = r_tx;
if (ray_opts.raytrace_frame != options_.frame) {
  r_rx_ray = ConvertFrame(t_tdb, r_rx, options_.frame, ray_opts.raytrace_frame);
  r_tx_ray = ConvertFrame(t_tdb, r_tx, options_.frame, ray_opts.raytrace_frame);
}

pecsim::PathProfile profile = pecsim::trace_ray(
    t_raytrace.val(), ToVec3d(r_tx_ray) / 1000.0, ToVec3d(r_rx_ray) / 1000.0, config,
    ray_opts.debug_propagation, ray_opts.debug_correction);

Real delay_m = profile.tec_delay_m;
if (ray_opts.delay_mode == GnssIonospherePlasmaRayTraceDelayMode::TEC_PLUS_HIGHER_ORDER)
  delay_m += profile.second_delay_m + profile.third_delay_m;
return delay_m;
\end{lstlisting}

The returned $\Delta\rho_\mathrm{plasma}$ is a positive code delay and enters the
observation equations of \cref{chap:gnss-measurements} with opposite signs on the
two observables,
%
\begin{equation}
  \label{eq:iono-observables}
  \rho = \ldots + \Delta\rho_\mathrm{plasma} + \ldots ,
  \qquad
  \varphi = \ldots - \Delta\rho_\mathrm{plasma} + \ldots ,
\end{equation}
%
reflecting the phase advance / group delay duality of \cref{sec:iono-signs}.

The \cpp{GnssIonospherePlasmaRayTraceOptions} fields are \code{config} (the
tracer knobs of \cref{tab:iono-config}; note that \cpp{GNSSMeasurements} never
fills in $K_p$, the integrator or the step sizes --- those must come from the
caller), \code{raytrace\_frame} (default \code{ECEF}),
\code{raytrace\_epoch\_scale} (default \code{UTC}),
\code{use\_channel\_frequency} (default true, which overwrites
\code{config.freq\_Hz} with the channel frequency), \code{iri\_model}, and
\code{delay\_mode}.

\begin{lupntnote}
\code{delay\_mode} defaults to \code{TEC\_ONLY}, so the delay returned to the
measurement model is \code{profile.tec\_delay\_m} alone, i.e.\
$d_{I1,\mathrm{LOS}} + d_{I1,B}$; \code{TEC\_PLUS\_HIGHER\_ORDER} adds
$d_{I2} + d_{I3}$.  The geometric excess-path term $d_\mathrm{len}$
(\code{profile.dist\_bend\_m}) is computed and reported but is \textbf{not} added
to the pseudorange: its timing effect is already captured by the
$\Delta\mathrm{TEC}_\mathrm{bend}$ part of \code{tec\_delay\_m}, and the residual
sub-metre geometric lengthening is deliberately left out of the range equation.
\end{lupntnote}

\section{Batch Precompute}
\label{sec:iono-precompute}

Ray tracing is expensive: one GCPM call per gradient component per stage per
step per correction iteration.  Two execution paths therefore exist.

The \emph{online} path calls \cpp{trace\_ray} link-by-link inside
\cpp{GNSSMeasurements::Compute}.  The \emph{batch} path decouples the tracer from
filtering: in the staged lunar GNSS ODTS scenario, \cpp{Precompute} writes all
light-time-corrected links with zero plasma delay, a Python driver fills each
epoch/channel delay by ray tracing on a coarse cadence and interpolating to
\SI{1}{\hertz}, and \cpp{Run} merges the delay file back before forming the
observables.  The two paths agree exactly on the measurement model; see the
online/precompute equivalence discussion in \cref{chap:gnss-measurements}.

\begin{lstlisting}[language=py, caption={Batch delay generation driver.  Source:
\texttt{python/examples/ex6\_precompute.py} (writes
\texttt{precomputed\_delays.csv}).}, label={lst:iono-precompute}]
# python/examples/ex6_precompute.py
FREQ_HZ = {"L1": 1575.42e6, "L2": 1227.60e6, "L5": 1176.45e6,
           "E1": 1575.42e6, "E5a": 1176.45e6}
EARTH_RADIUS_KM = 6371.2

def _rt_params(cfg):
    """Picklable dict of RayTraceConfig attribute values, sent to each worker."""
    p = cfg.plasma
    return dict(step_size=p.raytrace_step_size_km,
                correction=p.raytrace_correction,
                fine_correction=p.raytrace_fine_correction,
                cutoff_r=p.raytrace_cutoff_radius_re * EARTH_RADIUS_KM,
                gradn_dx=p.raytrace_gradient_step_km,
                integ_method=p.raytrace_integrator,
                correction_method=p.raytrace_correction_method,
                kp=p.raytrace_kp, rz12=p.raytrace_rz12,
                use_fortran_gcpm=p.raytrace_use_fortran_gcpm,
                corr_tol=p.raytrace_correction_tolerance_m,
                compute_higher_order=p.raytrace_compute_higher_order,
                use_adaptive_step=p.raytrace_use_adaptive_step,
                straight_ray=p.raytrace_straight_ray)

def _init_worker(params):
    """Runs in each worker (fork or spawn): select the IRI model and build its RayTraceConfig."""
    global _RT_CONFIG
    pnt.set_iri_model("IRI2007")
    _RT_CONFIG = _build_rt_config(params)

def _trace_worker(task):
    """Ray-trace one link through GCPM/IRI2007 -> dict of delay columns."""
    t_utc, tx_km, rx_km, freq_hz = task
    _RT_CONFIG.freq_Hz = freq_hz
    profile = pnt.trace_ray(t_utc, tx_km, rx_km, _RT_CONFIG, False, False)
    tec_m = float(profile.tec_delay_m)
    higher_m = (float(profile.second_delay_m) + float(profile.third_delay_m)
                if _RT_CONFIG.compute_higher_order else 0.0)
    return {"tecu": float(profile.tecu),
            "tec_delay_m": tec_m,
            "second_delay_m": float(profile.second_delay_m),
            "third_delay_m": float(profile.third_delay_m),
            "ionosphere_plasma_delay_m": tec_m + higher_m}
\end{lstlisting}

Two details of the driver are load-bearing.  First, \cpp{set\_iri\_model} mutates
process-global Fortran COMMON state and reads the IRI data files, so it must be
called once per worker process in the initializer rather than per ray.  Second,
the cached delays are keyed on a fingerprint that combines the Stage-1 link
fingerprint with every \cpp{RayTraceConfig} field and the interpolation cadence,
so changing $K_p$, $R_{12}$, the integrator or the cutoff invalidates the cache
while filter tuning does not.

A pure C++ batch entry point is also exposed through
\file{cpp/lupnt/environment/plasma/tec/raytrace_wrapper.h}
(\cpp{pecsim::run\_raytrace\_batch}, \cpp{load\_raytrace\_data},
\cpp{RaytraceData}), which stores per ray the fields listed in
\cref{tab:iono-batch}.  The $K_p$ time series consumed by these runs is fetched
from GFZ Potsdam by \file{python/pylupnt/plasma/kp_loader.py} and cached under
the plasma data directory.

\begin{table}[htbp]
  \centering
  \caption{Per-ray records in \cpp{RaytraceData}
  (\texttt{raytrace\_wrapper.h}).}
  \label{tab:iono-batch}
  \small
  \begin{tabularx}{\textwidth}{lL}
    \toprule
    Field & Contents \\
    \midrule
    \code{tidxs}, \code{row\_fulls}          & Transmitter index and full-table row index \\
    \code{epoch\_utcs}                       & Reception epoch [\si{\second}] \\
    \code{tx\_positions}, \code{rx\_positions} & Endpoint positions [\si{\meter}] \\
    \code{total\_delays}                     & $d_\mathrm{total}$ [\si{\meter}] \\
    \code{tecus}                             & Bent-path TEC [\si{\TECU}] \\
    \code{tec\_delays}                       & $d_{I1,\mathrm{LOS}} + d_{I1,B}$ [\si{\meter}] \\
    \code{second\_delays}, \code{third\_delays} & $d_{I2}$, $d_{I3}$ [\si{\meter}] \\
    \code{dist\_bend\_m}                     & $d_\mathrm{len}$ [\si{\meter}] \\
    \code{tec\_delay\_bend\_m}               & $d_{I1,B}$ [\si{\meter}] \\
    \code{min\_alts}                         & Minimum (tangential) altitude along the path [\si{\meter}] \\
    \code{max\_sep\_line\_m}                 & Maximum separation from the chord [\si{\meter}] \\
    \code{final\_pos\_err\_m}                & Residual terminal miss [\si{\meter}] \\
    \bottomrule
  \end{tabularx}
\end{table}

\section{Model Boundaries and Notes}
\label{sec:iono-boundaries}

\begin{description}[style=nextline, leftmargin=1.2em]
  \item[Km-unit convention.]
    All tracer inputs and outputs are km / km-s based; SI callers scale by
    $1000$ at the boundary.  Delays, however, are returned in metres.

  \item[Magnetic-field units.]
    The delay coefficients assume $B$ in \si{\nano\tesla}; the thesis symbolic
    coefficients are the SI-tesla equivalents (\cref{tab:iono-coeffs}).  The
    field components are additionally expressed in the local north/east/down
    triad while $\uvec{s}$ is geocentric Cartesian, so $\cos\theta$ is a proxy
    rather than an exact projection.

  \item[Second and third order are small.]
    For terrestrial and LEO users they stay below about \SI{1}{\percent} of the
    first-order delay; for the lunar geometry
    \citep[Tab.~5.2]{iiyama2026thesis} places them several orders of magnitude
    below $d_{I1}$.  The GNSS default (\code{TEC\_ONLY}) omits them, and they are
    the only consumers of the IGRF evaluation.

  \item[Cutoff at $4\,R_E$.]
    Plasma outside the inner magnetosphere is neglected entirely and the path is
    taken to be exactly straight beyond the cutoff.  For a lunar receiver at
    $\approx 60\,R_E$ this means roughly \SI{93}{\percent} of the geometric path
    length contributes nothing to the delay --- which is correct to well within
    the model error, since the magnetosheath and solar-wind densities are
    $\lesssim \SI{10}{\per\cubic\centi\meter}$.

  \item[Time integration.]
    $\dd t/\dd s$ uses the vacuum $1/c$ with the plasma delay added separately;
    this is algebraically equivalent to $\dd t/\dd s = n/c$ but avoids
    catastrophic cancellation against an epoch of order $10^{8}\,\si{\second}$.

  \item[Excess path length.]
    $d_\mathrm{len}$ is characterized and reported but is \emph{not} added to the
    measurement delay (\cref{sec:iono-gnss}).

  \item[Geometric optics only.]
    The eikonal formulation assumes the WKB regime: no diffraction, no
    scintillation, no mode coupling between the ordinary and extraordinary
    branches of \cref{eq:iono-ah}, and no Faraday rotation bookkeeping.  Small
    scale irregularities below the \SIrange{10}{100}{\kilo\meter} step are not
    resolved, and GCPM does not contain them in the first place.

  \item[Empirical-model validity.]
    GCPM and IRI require historical solar and geomagnetic inputs and silently
    fall back to earlier years beyond their data coverage.  The characterization
    covers $K_p \in [1,9]$ and $R_{12} \in [10,200]$.

  \item[Magnitudes.]
    The delay is strongly dependent on the tangential (minimum) altitude of the
    path: the mean first-order delay is $\approx \SI{50}{\meter}$ for tangential
    altitudes below \SI{500}{\kilo\meter}, and falls below
    $\approx \SI{0.2}{\meter}$ above \SI{10000}{\kilo\meter}
    \citep{iiyama2026gcpm}.  Screening links by minimum altitude
    (\code{RaytraceData::min\_alts}) is therefore an effective, cheap
    approximation to running the full tracer on every link.
\end{description}
